\section{Indledning}
I anledning af, at Viborg Katedralskole fejrer bygningens 100-års jubilæum,
valgte gruppen at udvikle et spil, der både fejrer og hylder skolens historie,
arkitektur og kulturelle betydning. Gruppen ønskede at skabe et projekt,
som kombinerer moderne teknologi med historisk formidling, og
som samtidig kunne give brugeren en interaktiv oplevelse af skolens omgivelser og historie.
Motivation for projektet opstod ud fra et ønske om at arbejde med et lokalt og meningsfuldt emne,
hvor der kunne skabes en tæt forbindelse mellem det tekniske produkt og den historiske fortælling.
\\\\
I projektets opstartsfase arbejdede gruppen med idéudvikling og research omkring skolens historie
og arkitektoniske udtryk. Gennem møder med skolens historielærere og gennemgang af historisk
materiale blev gruppen introduceret til skolens rolle under Anden Verdenskrig.
Her blev det blandt andet beskrevet,
hvordan modstandsfolk benyttede skolen til produktion og distribution af illegale aviser
under den tyske besættelse. Gruppen fandt denne fortælling særligt interessant,
da den både havde historisk betydning og samtidig kunne omsættes til spændende gameplay-elementer
 et digitalt spil.
\\\\
På baggrund af dette valgte gruppen at udvikle et online multiplayer-spil,
hvor spillerne enten kan spille som modstandsfolk eller som nazister.
Modstandsfolkene skal bevæge sig rundt på skolen og udføre forskellige opgaver,
såsom at Finde Nøgler som refere til Skolens Logo og Historie eller undgå at blive opdaget,
mens den modsatte side forsøger at finde og stoppe dem.
Gruppen ønskede gennem dette gameplay at skabe spænding og samarbejde mellem spillerne,
samtidig med at spillet tog inspiration fra virkelige historiske begivenheder.
\\\\
En vigtig del af projektet var at skabe et visuelt realistisk miljø,
som kunne give spillerne følelsen af faktisk at befinde sig på skolen.
Derfor blev spillet udviklet i Unreal Engine, da gruppen vurderede,
at denne engine gav de bedste muligheder for høj grafisk kvalitet,
realistisk lysopsætning og detaljerede omgivelser.
Gruppen havde desuden et ønske om at lære mere om professionel spiludvikling og arbejde med værktøjer,
der anvendes i den moderne spilindustri.
\\\\
For at gøre omgivelserne så virkelighedstro som muligt anvendte gruppen 3D-scanninger
af skolens bygninger og områder. Dette blev gjort ved hjælp af programmet RealityScan,
som gjorde det muligt at omdanne billeder af bygningerne til detaljerede 3D-modeller såkaldt photogrammetry.
Disse modeller blev efterfølgende importeret til spillet og viderebearbejdet,
så de kunne anvendes i det digitale miljø.
Gruppen oplevede, at dette både gav en mere autentisk oplevelse
og samtidig gjorde projektet mere teknisk udfordrende,
da arbejdet med 3D-scanning og optimering krævede nye kompetencer
inden for modellering og performanceoptimering.
\\\\
Projektforløbet blev organiseret ud fra en Scrum-lignende arbejdsmetode med sprintbaseret udvikling.
Gruppen valgte denne agile tilgang for at skabe en fleksibel arbejdsproces,
hvor der løbende kunne arbejdes med planlægning, udvikling, test og evaluering.
Hvert sprint havde fokus på bestemte funktioner eller dele af spillet,
hvilket gjorde det muligt at arbejde mere struktureret og samtidig tilpasse projektet undervejs.
Gruppen anvendte blandt andet møder og opgavefordeling,
for at sikre overblik over arbejdsprocessen og kommunikationen mellem gruppemedlemmerne.
\\\\
Igennem projektet oplevede gruppen både tekniske og samarbejdsmæssige udfordringer.
Særligt arbejdet med multiplayer-funktioner,
3D-modeller og koordinering mellem forskellige arbejdsområder krævede løbende problemløsning
og samarbejde.
Samtidig gav projektet gruppen mulighed for at udvikle kompetencer inden for programmering,
projektstyring, kommunikation og kreativ problemløsning.
Gruppen fik desuden erfaring med,
hvordan historiske fortællinger kan integreres i digitale medier for at skabe en engagerende brugeroplevelse.
\\\\
Procesrapporten vil derfor beskrive projektets udviklingsforløb,
arbejdsmetoder og samarbejde samt gruppens refleksioner over læringsudbytte,
udfordringer og erfaringer gennem hele projektperioden.